\secciongrande{Herramientas de programación sobre clientes web: Elección del editor de código}

\resumen{
  La elección de una herramienta de desarrollo depende del programador, pudiendo ser un editor de texto simple. Es crucial seleccionar desarrolladores de código y no editores de contenido como FrontPage. Las pautas incluyen soporte para codificación de colores, verificación de sintaxis y opciones como clientes FTP. Editores gratuitos recomendados son Notepad++, Eclipse, NetBeans, KompoZer y Bluefish.
}

\definicion{Editor de código}{Herramienta especializada en la escritura, edición y gestión de código fuente, optimizada para lenguajes de programación, a diferencia de los editores de contenido orientados a la presentación.}
\subtituloA{Pautas para la selección de herramientas}

\lista{
  \textbf{No utilizar editores de contenido:} Evitar herramientas como FrontPage o Dreamweaver, enfocándose en desarrolladores de código. |
  \textbf{Facilidad para desarrollo web:} Buscar codificación de colores, verificación de sintaxis, inserción de comentarios, generación de blogs y clientes FTP.
}
\nota{Para el examen}{
Diferenciar entre un editor de código (enfocado en la sintaxis) y un editor de contenido (enfocado en la presentación visual).
}
